\subsection{The complete compact-unitary Hecke models in three further weights}
\label{sec:u3-general-harmonics}

\begin{theorem}[K17-U3-GENERAL]\label{thm:u3-general}
Retain the standard positive Hermitian space $E^3$, with
$E=\mathbf Q(\omega)$, $\omega^2+\omega+1=0$, and its full integral
lattice stabilizer $K$ from Theorem~\ref{thm:u3-hecke}.
For $e=8,10,12$, let $\mathcal H_{e-1}^{\Gamma}$ be the entire space
of algebraic modular forms of highest weight $(e-1,0,1-e)$ at level $K$.
Its dimensions are respectively $3,6,8$.
For $p=7,13$, let $T_p^+$ be the unnormalized positive minuscule
Hecke operator, with compact-subgroup volume one, and put
\[
 P_{e,p}(X)=\det\bigl(X-p^{e-1}T_p^+\mid
                       \mathcal H_{e-1}^{\Gamma}\bigr).
 \tag{UG.1}
\]
These are integral polynomials. The complete exact factorizations are
given in (UG.10)--(UG.12) below. For $E_{e,p}=1+p^e+p^{2e}$ one has
\[
\begin{array}{c|l}
e&\gcd\bigl(P_{e,7}(E_{e,7}),P_{e,13}(E_{e,13})\bigr)\\ \hline
8&2^{13}3^{13}\,23\,401\,809\,13687\\
10&2^{28}3^{26}\,43\,241\,1847\,9496386083\\
12&2^{28}3^{39}\,13^4\,47\,28447\,300455201.
\end{array}
\tag{UG.2}
\]
Every displayed prime factor is certified by trial division. The
operator in (UG.1) has eigenvalues equal to $p^e$ times the standard
normalized Satake trace. This theorem concerns the specified full
coefficient spaces and full lattice level; the arithmetic source
interface is a separate assertion.
\end{theorem}

\begin{proof}
The explicit Hermitian reduction in Theorem~\ref{thm:u3-hecke} proves
class number one without a weight restriction. Its integral stabilizer
is $\Gamma=\mu_6^3\rtimes S_3$. Evaluation at the sole adelic class
therefore identifies the whole space with the $\Gamma$-invariants in
the indicated irreducible representation. We first construct that
representation and every one of its invariants, then identify the full
Hecke matrices by injective exact evaluation.

\paragraph{The whole harmonic representation.}
Use independent formal variables $z=(z_1,z_2,z_3)$ and
$w=(w_1,w_2,w_3)$, and put
\[
 R=\sum_i z_iw_i,\qquad
 \Delta=\sum_i\partial_{z_i}\partial_{w_i},\qquad
 \rho(g)f(z,w)=f(g^{-1}z,\overline g^{-1}w).
 \tag{UG.3}
\]
This is the complexification of the usual polynomial model in $z$ and
$\overline z$. Let $\mathcal H_k$ be the kernel of $\Delta$ in
bidegree $(k,k)$. The contraction onto bidegree $(k-1,k-1)$ is
surjective: for $F$ of that bidegree, a preimage is
\[
 \sum_{j=0}^{k-1}c_jR^{j+1}\Delta^jF,\qquad
 c_0=\frac1{2k+1},\qquad
 c_j=-\frac{c_{j-1}}{(j+1)(2k+1-j)}.
 \tag{UG.4}
\]
Indeed, the product rule gives
\[
 \Delta(R^{j+1}\Delta^jF)
 =(j+1)(2k+1-j)R^j\Delta^jF+R^{j+1}\Delta^{j+1}F,
\]
and consecutive terms cancel. Hence
\[
 \dim\mathcal H_k
 =\binom{k+2}{2}^{\!2}-\binom{k+1}{2}^{\!2}=(k+1)^3.
\]
It contains the highest-weight vector $w_1^kz_3^k$ of weight
$(k,0,-k)$, whose irreducible representation has this same Weyl
dimension. Complete reducibility for the compact group identifies
the entire kernel with that irreducible representation.

\paragraph{Every finite-group invariant.}
An invariant monomial $z^\alpha w^\beta$ must satisfy
\[
 |\alpha|=|\beta|=k,\qquad
 \alpha_i-\beta_i\equiv0\pmod6\quad(1\le i\le3).
 \tag{UG.5}
\]
For each simultaneous $S_3$-orbit of the three pairs
$(\alpha_i,\beta_i)$, take the sum of its distinct monomials.
These orbit sums have disjoint supports and form a basis of all
$\Gamma$-invariants. Thus (UG.5) includes the additional exponent
differences $\pm6$ in the present degrees. It does not impose the
incorrect condition $\alpha=\beta$ once $k\ge6$.
The complete enumeration consists of all pairs of three-part
compositions of $k$; for $k=11$ there are only $78^2$ such pairs.

The contraction matrix follows from
\[
 \Delta(z^\alpha w^\beta)
 =\sum_i\alpha_i\beta_i z^{\alpha-e_i}w^{\beta-e_i}.
\]
The discovery program collects coefficients by target orbit and
divides by its size. The independent certificate expands every
monomial, reads the coefficient of a canonical target monomial,
and checks equality on every monomial in that target orbit.
The resulting exact ranks are
\[
\begin{array}{c|c|r|r|r|c}
k&e&\dim V_k^\Gamma&\dim V_{k-1}^\Gamma&
 \dim\mathcal H_k^\Gamma&\text{swap dimensions }(+,-)\\ \hline
5&6&5&4&1&(1,0)\\
7&8&11&8&3&(2,1)\\
9&10&22&16&6&(4,2)\\
11&12&37&29&8&(5,3).
\end{array}
\tag{UG.6}
\]
Here $V_k$ denotes all polynomials of bidegree $(k,k)$; $k=5$ is the
previously proved control. Each saved harmonic basis consists of
integral columns annihilated by the full contraction matrix.
Their ranks and the contraction ranks establish that they span
the complete characteristic-zero kernel. No saturation assertion
for this particular integral basis is needed.

\paragraph{Actual cosets, action, and direction.}
At each split prime $p$, choose $\pi\in\mathbf Z[\omega]$ of norm $p$.
The complete negative-minuscule representatives in the certificate
are $g=U/\pi$ with
$U^*U=pI$ and $\operatorname{rank}(U\bmod\pi)=1$.
There is one for each line of $\mathbf P^2(\mathbf F_p)$, hence
$57$ for $p=7$ and $183$ for $p=13$.
The earlier lattice and elementary-divisor argument identifies
these lines with the actual right cosets: the elementary divisors
of $U$ at $\pi$ are $(0,1,1)$, so those of $g$ are $(-1,0,0)$;
the conjugate lattice is forced by Hermitian duality, and the
lattice is unchanged at other primes. Every matrix identity and
every reduction line is checked again by the present certificate,
as is exact coverage of the entire projective plane.

Since $g^{-1}=U^*/\overline\pi$ and
$\overline g^{-1}=U^t/\pi$, homogeneity gives the exact formula
\[
 p^k\rho(g)f(z,w)=f(U^*z,U^tw).
 \tag{UG.7}
\]
The Hecke operator is the sum over the complete coset set. It
preserves the whole invariant harmonic space: the Hecke
correspondence preserves $\Gamma$-invariance and unitary
substitution preserves harmonicity.

For clarity, positive and negative directions are compared on
the entire space. Let $C$ interchange $z$ and $w$. Then
$C\rho(g)C=\rho(\overline g)$, while conjugation bijects the
negative and positive minuscule coset sets. Thus their matrices
are conjugate by $C$ and their full characteristic polynomials
agree. This argument does not require either operator to be scalar.
Choose the harmonic basis in the $C$-eigenspaces and write
$C=\operatorname{diag}(\epsilon_i)$ and the negative scaled
matrix as $M=M_0+\omega M_1$. Coefficientwise conjugation is the
positive operator, so
\[
 CM_0C=M_0-M_1,\qquad CM_1C=-M_1.
 \tag{UG.8}
\]
The certificate checks these identities entry by entry.
Put $\delta=1+2\omega$, so $\delta^2=-3$, and multiply each
$C$-minus basis vector by $\delta$. The resulting matrix is
rational. Same-sign entries have $M_1=0$ and remain $M_0$;
opposite-sign entries have $M_1=2M_0$ and become $-3M_0$ in
a plus row and minus column, or $M_0$ in a minus row and plus
column. All these rational entries are recorded exactly.

The characteristic polynomial is integral. Formula (UG.7) shows
that the scaled operator preserves the lattice of invariant
harmonic polynomials with coefficients in $\mathbf Z[\omega]$.
This is a full finite free lattice over that PID. The
characteristic polynomial therefore lies in $\mathbf Z[\omega][X]$;
(UG.8) places it in $\mathbf Q[X]$, and the intersection is
$\mathbf Z[X]$.
Finally, Gross's minuscule Satake formula
\cite[\S3, (3.13)--(3.14), p.~8]{K17GrossSatake1998}
identifies $T_p^+$ with $p$ times the standard normalized trace.
Thus the scaling $p^k=p^{e-1}$ in (UG.1) gives $p^e$ times that
trace. There is no division by $|\Gamma|$ or extra stabilizer factor.

\paragraph{Injective exact evaluation identifies the whole matrix.}
Let $h_1,\ldots,h_d$ be the saved full harmonic basis. Choose $d$
pairs $(z_r,w_r)$ of integer triples and set
\[
 A_{rj}=h_j(z_r,w_r),\qquad
 B_{rj}=\sum_U h_j(U^*z_r,U^tw_r)=B_{0,rj}+\omega B_{1,rj}.
\]
Every point and every entry of $A$ is saved, and its exact
determinant is nonzero. These are evaluations of the complexified
polynomials in independent variables, so no restriction to the
real locus $w=\overline z$ is appropriate or necessary.
Because the scaled operator already preserves the whole space,
the equations
\[
 AM_0=B_0,\qquad AM_1=B_1
 \tag{UG.9}
\]
determine its entire matrix. The independent standard-library
program reconstructs $A$, verifies its determinant, recomputes
every $B$ entry with Eisenstein integer pairs, and verifies
(UG.9) over exact fractions. Two further point pairs per prime
check the same operator identity; they test the implementation
and are not substitutes for the preceding completeness argument.

The certificate next computes the characteristic polynomial by
Faddeev--LeVerrier over exact fractions, verifies the terminal
Cayley--Hamilton zero matrix, and verifies that all coefficients
are integers. It multiplies the factors below and compares them
with that complete polynomial. No irreducibility of those factors
is needed or asserted.

\paragraph{The exact polynomial outputs.}
For $e=8$ the factors are
\[
\begin{aligned}
 P_{8,7}(X)&=(X+2702301)
             (X^2-9958830X-15191091292671),\\
 P_{8,13}(X)&=(X+153697821)
             (X^2-841096230X+74635873683999849).
\end{aligned}
\tag{UG.10}
\]
For $e=10$, write $P_{10,p}=L_{10,p}A_{10,p}B_{10,p}$, where
\[
\begin{aligned}
 L_{10,7}&=X-14535507,\\
 A_{10,7}&=X^2+571295802X+259067515101993801,\\
 B_{10,7}&=X^3-1127229705X^2+269867622745258587X\\
          &\quad+21764201659661735207968293,\\
 L_{10,13}&=X+109076012133,\\
 A_{10,13}&=X^2+106910299722X+2996591715398898111321,\\
 B_{10,13}&=X^3-422499415185X^2+52678214627106136408587X\\
          &\quad-1658111541578384339481222489985827.
\end{aligned}
\tag{UG.11}
\]
For $e=12$, write $P_{12,p}=L_{12,p}A_{12,p}B_{12,p}$, where
\[
\begin{aligned}
 L_{12,7}&=X+12221224701,\\
 A_{12,7}&=X^3-42517894041X^2+288396328713382721979X\\
          &\quad-463507446167295652333685832939,\\
 B_{12,7}&=X^4+2939150868X^3-79381481904803039850X^2\\
          &\quad+303898318630502425871984942580X\\
          &\quad+8681328863742961883955067744457913482097,\\
 L_{12,13}&=X+21649300719981,\\
 A_{12,13}&=X^3-33123036335481X^2-228653307502980173933740101X\\
          &\quad+3302349716440899212738561798170057697781,\\
 B_{12,13}&=X^4+48903137840436X^3+2568930841937762809121509110X^2\\
          &\quad+39746045119592037993985794601727262212436X\\
          &\quad+940042707779377699684081382859793759142653810554146961.
\end{aligned}
\tag{UG.12}
\]
The complete expanded coefficients, all rational matrices, and
every identifying evaluation are in the frozen data file.
Integer Horner evaluation at $1+p^e+p^{2e}$ and the Euclidean
algorithm give (UG.2). The verifier multiplies every listed
factor and proves its primality by trial division through its
integer square root; the largest bound is below $100000$.
This completes the exact finite calculation.
\end{proof}

\paragraph{Reproduction and the prime nineteen check.}
The complete data and standard-library verifier are
\path{hecke_general/harmonic_hecke_exact_data.json} and
\path{hecke_general/harmonic_hecke_certificate.py}.
The frozen data and the complete neighbor input are pinned by
their SHA256 values in the replay report. Every neighbor is
independently checked again as part of that replay.
The discovery program uses symbolic rational linear algebra;
the replay imports only the Python standard library and does not
call or trust that symbolic solve. Both have been run successfully
on all the displayed degrees and primes. The saved replay records
all basis and contraction ranks, every neighbor check, every
transformed evaluation, the full characteristic polynomials,
and the common-divisor factorizations.

The same verifier, with explicitly selected alternate input files
under \texttt{hecke\_general/p19/}, also checks the $381$ neighbors
at $p=19$, using $\pi=5+2\omega$, and the full six-dimensional
space for $e=10$. Its polynomial is
\[
\begin{aligned}
 P_{10,19}(X)&=(X+5229588168357)\,A(X)\,B(X),\\
 A(X)&=X^2-10082283620118X+30950864549521920787631481,\\
 B(X)&=X^3-10914188675841X^2-3704020704004235185284885X\\
     &\quad+53864419616244584992508537884872333357.
\end{aligned}
\tag{UG.13}
\]
The short standard-library file
\texttt{p19/three\_prime\_comparison\_certificate.py} evaluates all
three full polynomials and verifies that their common divisor is
\[
 2^{28}3^{26}\,241\,1847\,9496386083.
 \tag{UG.14}
\]
Thus the extra test removes the factor $43$ from the U3 common divisor
and retains $9496386083$. It does not eliminate the latter prime.
An application using an auxiliary prime $p$ in residue characteristic
$\ell$ must retain $\ell\ne p$. Neither the two-prime nor the
three-prime calculation alone asserts new Vandiver verification,
global packet coverage, Selmer vanishing, or a new algebraic $K$-group
result.
